\documentclass[10pt]{article}
\usepackage[accepted]{icml2026}

\usepackage[T1]{fontenc}
\usepackage[utf8]{inputenc}
\usepackage{lmodern}
\usepackage{microtype}
\usepackage{amsmath,amssymb}
\usepackage{booktabs}
\usepackage{enumitem}
\usepackage[numbers,sort&compress]{natbib}
\usepackage{xcolor}
\usepackage{hyperref}
\usepackage{float}

\hypersetup{
  colorlinks=true,
  linkcolor=blue!50!black,
  citecolor=blue!50!black,
  urlcolor=blue!50!black
}

\newcommand{\RESULTPENDING}{\textcolor{red!70!black}{\textbf{RESULT PENDING}}}

\icmltitlerunning{PolyBench}

\title{\textbf{PolyBench: A Real-Time Prediction Market Benchmark for AI Forecasting Agents}}
\author{}
\date{}

\begin{document}

\twocolumn[
\maketitle
\begin{center}
\vspace{-1.75em}
\begin{minipage}{0.92\textwidth}
\begin{abstract}
Current AI forecasting benchmarks mostly evaluate forecasters on static question
sets with delayed feedback and no continuously updating market baseline.
Prediction markets offer a different regime: probabilities update in real time,
settlement is objective, and the benchmark includes a strong crowd-implied
reference forecast at every evaluation point. We propose PolyBench, a benchmark
for forecasting agents built on a live Polymarket data stack. The underlying
repository already ingests market metadata, trades, orderbooks, wallet activity,
and derived analytics into ClickHouse; it also contains calibration tracking,
backtesting, and signal infrastructure that can be repurposed for agent
evaluation. PolyBench evaluates an agent not only by proper scoring rules such
as Brier score and log score, but also by \emph{informational alpha}: the
marginal value of an agent relative to the market-implied baseline at the moment
the forecast is made. The current draft defines the benchmark task, scoring
protocol, and minimum viable implementation path for a workshop submission.
\end{abstract}
\end{minipage}
\end{center}
\vspace{0.6em}
]

\section{Introduction}

Forecasting benchmarks have improved quickly, but they still leave an important
gap between forecasting as a question-answering task and forecasting as a live
decision problem. ForecastBench is valuable because it is dynamic and
leakage-resistant, but it does not provide a continuously updating market
baseline \citep{karger2024forecastbench}. PredictionMarketBench moves closer to
market realism by replaying exchange data for trading agents, but it focuses on
execution-realistic backtesting rather than on a general benchmark for
probability forecasters operating against a live crowd baseline
\citep{arora2026predictionmarketbench}.

PolyBench targets the missing middle:
\begin{itemize}[leftmargin=1.5em,itemsep=0.2em]
  \item a real-time benchmark;
  \item objective resolution;
  \item a strong crowd-implied baseline at every timestamp;
  \item difficulty stratification by liquidity, spread, and horizon;
  \item evaluation that rewards \emph{useful} disagreement with the market.
\end{itemize}

\section{Benchmark Task}

At time $t$, an agent observes a benchmark market and emits a probability
forecast $p_{\text{agent}}(t)$ for the event's positive outcome. For each
forecast event, PolyBench records:
\begin{itemize}[leftmargin=1.5em,itemsep=0.2em]
  \item market identifier and timestamp;
  \item market-implied probability at that timestamp;
  \item spread, depth, and other liquidity-quality features when available;
  \item category and horizon metadata;
  \item the agent forecast;
  \item the eventual resolved outcome.
\end{itemize}

The benchmark is designed to support two modes:
\begin{itemize}[leftmargin=1.5em,itemsep=0.2em]
  \item \textbf{live mode}, where agents forecast active markets;
  \item \textbf{replay mode}, where historical market states are served
  deterministically for reproducible evaluation.
\end{itemize}

\section{Scoring}

PolyBench uses standard proper scoring rules plus a differential metric that
explicitly measures value added over the market.

\paragraph{Core metrics.}
We report Brier score, log score, Expected Calibration Error, and slice-level
coverage across domain and horizon \citep{brier1950verification}.

\paragraph{Informational alpha.}
For forecast instance $i$, define
\[
\alpha_i = (p_{\text{market},i} - y_i)^2 - (p_{\text{agent},i} - y_i)^2,
\]
where $y_i$ is the realized binary outcome. Positive alpha means the agent beat
the market on that forecast; negative alpha means it added noise rather than
information. The benchmark should report:
\begin{itemize}[leftmargin=1.5em,itemsep=0.2em]
  \item mean informational alpha;
  \item disagreement-conditioned alpha for forecasts with
  $|p_{\text{agent}} - p_{\text{market}}| \ge \tau$;
  \item liquidity-stratified alpha;
  \item category-stratified alpha.
\end{itemize}

\section{Why The Repo Makes This Feasible}

The local \texttt{polymarket} repository already contains most of the data
substrate needed for a benchmark:
\begin{itemize}[leftmargin=1.5em,itemsep=0.2em]
  \item continuous ingestion of market metadata, trades, orderbooks, and
  WebSocket events;
  \item derived tables for OHLCV, latest prices, wallet activity, arbitrage,
  composite signals, and market similarity;
  \item calibration tracking code in
  \texttt{pipeline/bayesian/calibration.py};
  \item backtesting and calibration analysis code in
  \texttt{network/backtest.py}.
\end{itemize}

What does \emph{not} yet exist as a polished benchmark artifact is just as
important:
\begin{itemize}[leftmargin=1.5em,itemsep=0.2em]
  \item a stable benchmark episode specification;
  \item a public agent interface;
  \item dedicated baseline-agent runners;
  \item a canonical replay service.
\end{itemize}
This means the workshop paper should be framed as a systems and benchmark paper
with a minimum viable implementation rather than as a mature public benchmark.

\section{Baseline Agents}

The first submission needs only a small set of transparent baselines.

\paragraph{Market Echo.}
Always predict the current market-implied probability. This is the zero-alpha
reference point and the benchmark baseline to beat.

\paragraph{Momentum.}
Use recent market movement and activity features to shift the forecast away from
the market when short-term price dynamics imply continued drift. The current
repo already exposes candidate inputs for such a baseline through one-day price
change, trade imbalance, and signal-compositor outputs.

\paragraph{LLM-RAG forecaster.}
This is the most interesting long-run baseline, but it is not required for the
first paper. Proposal B remains valuable even if the initial submission only
contains Market Echo and Momentum plus the benchmark specification.

\section{Current Status}

The benchmark framing is concrete, but the benchmark-specific empirics are still
behind Proposal A:
\begin{itemize}[leftmargin=1.5em,itemsep=0.2em]
  \item replay episode schema: \RESULTPENDING
  \item Market Echo baseline run: \RESULTPENDING
  \item Momentum baseline run: \RESULTPENDING
  \item category- and liquidity-sliced alpha table: \RESULTPENDING
\end{itemize}

The right sequencing is therefore to use Proposal A as the primary empirical
submission while building Proposal B in parallel as the systems companion.

\section{Conclusion}

Prediction markets enable a benchmark regime that standard forecasting datasets
do not: agents can be evaluated against a continuously updating crowd baseline
with objective settlement and natural difficulty variation across liquidity,
horizon, and category. The local Polymarket stack already contains most of the
required data machinery. The main remaining work is packaging that machinery
into a clean replay benchmark with a transparent baseline suite and a scoring
protocol centered on informational alpha.

\bibliographystyle{plainnat}
\bibliography{references}

\end{document}
